%%%%%%%%%%%%%%%%%%%%%%%%%%%%%%%%%%%%%%%%%%%%%%%%%%%%%%%%%%%%%%%%%%%%%%%%%%%%%%%
% Lossless compression
%%%%%%%%%%%%%%%%%%%%%%%%%%%%%%%%%%%%%%%%%%%%%%%%%%%%%%%%%%%%%%%%%%%%%%%%%%%%%%%

\subsection{The Source Coding Problem}
\label{sec:lossless-intro}

\begin{figure}[h]
    \centering
    \includegraphics[width=0.78\textwidth]{scribe3-SC.png}
    \caption{Lossless source coding: an encoder maps symbols to bit strings, and the decoder must recover the symbols with zero error.}
    \label{fig:source-coding}
\end{figure}

In \emph{lossless} compression the goal is to represent data exactly (no distortion), using as few bits as possible.
We start with the simplest model: a discrete memoryless source.

\paragraph{Discrete memoryless source.}
Let $X$ be a random variable on a finite alphabet $\cX$ with pmf $p_X$.
The source produces a sequence $X_1,X_2,\dots$ of i.i.d. draws, each distributed as $X$.

\subsection{Variable-Length Codes and Rate}

A (binary) \emph{code} assigns each symbol $x\in\cX$ a finite bitstring (a \emph{codeword})
\(f(x)\in\{0,1\}^*\).
Let $\ell(x)\coloneqq |f(x)|$ denote its length.
Encoding a sequence $x^n=(x_1,\dots,x_n)$ produces the concatenation
\(f(x_1)f(x_2)\cdots f(x_n)\).

\begin{definition}[Expected rate]
For a symbol-by-symbol code, the expected number of bits per symbol (the \emph{rate}) is
\begin{equation}
    L \coloneqq \E{\ell(X)} = \sum_{x\in\cX} p_X(x)\,\ell(x). 
\end{equation}
\end{definition}

\subsection{Unique Decodability and Prefix-Free Codes}

\begin{definition}[Uniquely decodable]
A code $f$ is \emph{uniquely decodable} if every finite bitstring that can be produced by concatenation of codewords has a unique parsing into codewords.
Equivalently, $f(x^n)=f(\tilde x^m)$ implies $n=m$ and $x^n=\tilde x^m$.
\end{definition}

\begin{definition}[Prefix-free]
A code is \emph{prefix-free} if no codeword is a prefix of another.
Prefix-free codes are instantaneously decodable (decoding can proceed symbol-by-symbol from left to right).
\end{definition}

\begin{example}
Let $\cX=\{A,B,C\}$.
The code $A\mapsto 0$, $B\mapsto 10$, $C\mapsto 11$ is prefix-free.
In contrast, $A\mapsto 0$, $B\mapsto 1$, $C\mapsto 10$ is \emph{not} uniquely decodable because the bitstring $10$ could be $B$ or $CA$.
\end{example}

\begin{figure}[h]
    \centering
    \includegraphics[width=0.42\textwidth]{prefixcode_example.png}
    \caption{A prefix-free code corresponds to leaves of a binary tree; no leaf is an ancestor of another.}
    \label{fig:prefix-tree}
\end{figure}

\subsection{Kraft and Kraft--McMillan Inequalities}
\label{sec:kraft}

The central combinatorial constraint on codeword lengths is the Kraft inequality.

\begin{theorem}[Kraft inequality]
\label{thm:kraft}
Given positive integers $\{\ell(x):x\in\cX\}$, there exists a \emph{binary} prefix-free code with these codeword lengths if and only if
\begin{equation}
    \sum_{x\in\cX} 2^{-\ell(x)} \le 1.
\end{equation}
\end{theorem}

\begin{proof}
Suppose that $l_1 \leq l_2 \leq \cdots \leq l_n$. Let $A$ be a full binary tree of depth $l_n$. Each word of length $l_i$ corresponds to a node $v_i$ in this tree at length $l_i$. The subtree $A_i$ of $A$ rooted at $v_i$ has $2^{l_n-l_i}$ elements. And by the prefix-free property, $A_i\cap A_j=\emptyset\text{ for }i\neq j.$  Thus, $$
\Biggl|\bigcup_{i=1}^nA_i\Biggr| =\sum_{i=1}^n2^{l_n-l_i}\leq2^{l_n}=|A|.
$$

Conversely, set of codeword lengths $\{l_i\,:\,i=1,...,n\}$ satisfying this inequality, one can construct a prefix-free code by a greedy procedure.
\end{proof}


\begin{theorem}[Kraft--McMillan inequality]
\label{thm:kraft-mcmillan}
If a binary code is uniquely decodable with lengths $\{\ell(x)\}$, then
\(\sum_x 2^{-\ell(x)}\le 1\).
\end{theorem}


\begin{proof}
Denote $T=\sum_{i=1}^n2^{-l_i}.$ The idea of the proof is to get an upper bound on $T^n$ and show that it can only hold for all $n\in\mathbb{N}$ if $T \leq 1$.
Rewrite $T^n$ as
$$
\begin{aligned}
    T^n = \left(\sum_{i=1}^n 2^{-l_i}\right)^n=\sum_{i_1=1}^n\sum_{i_2=1}^n\cdots\sum_{i_n=1}^n2^{-(l_{i_1}+l_{i_2}+\cdots +l_{i_n})}.
\end{aligned}
$$
Consider $\mathcal{X}^n$. Since $\mathcal{X}$ was assumed to uniquely decodable, we can rewrite the equation to
$$T^n = \sum_{l=1}^{n\cdot l_{\max}} A_l 2^{-l}$$,
where $A_l$ corresponds to number of $l$-length words. Since $A_l\leq 2^{l}$, we can finally obtain
$$T^n = \sum_{l=1}^{n\cdot l_{\max}} A_l 2^{-l}\leq n\cdot l_{\max} \implies T=\sum_{i=1}^n 2^{-l_i}\leq (n \cdot l_{\max})^{\frac{1}{n}}$$. The right-hand side is $1$ asymptotically.
\end{proof}


\begin{remark}
Kraft--McMillan says that \emph{every} uniquely decodable code has lengths that satisfy the Kraft inequality, even if it is not prefix-free.
Moreover, for any uniquely decodable code there exists a prefix-free code with the same lengths.
So, when optimizing expected length, we may restrict attention to prefix-free codes.
\end{remark}

\subsection{Entropy Bounds for Lossless Coding}
\label{sec:entropy-bound}

Entropy controls the best achievable rate.

\begin{theorem}[Entropy lower bound]
\label{thm:entropy-lb}
For any prefix-free binary code for $X$,
\begin{equation}
    \E{\ell(X)} \ge H(X).
\end{equation}
\end{theorem}

\begin{proof}
Let $\ell(x)$ be the code lengths.
Define
\(Q(x)\coloneqq 2^{-\ell(x)}/K\), where $K\coloneqq\sum_{x}2^{-\ell(x)}\le 1$ by Kraft.
Then $Q$ is a pmf on $\cX$.
Using the cross entropy decomposition (Theorem~\ref{thm:cross-entropy-decomp}),
\[
\E{\log\frac{1}{Q(X)}}=H(X)+\KL{p_X}{Q}\ge H(X).
\]
On the other hand,
\[
\log\frac{1}{Q(x)} = \ell(x) + \log K^{-1} \le \ell(x)
\qquad (\text{since }K\le 1).
\]
Taking expectations gives $\E{\ell(X)}\ge \E{\log(1/Q(X))}\ge H(X)$.
\end{proof}

\begin{theorem}[Shannon code upper bound]
\label{thm:shannon-code}
There exists a prefix-free code (a \emph{Shannon code}) with lengths
\(\ell(x)=\lceil \log (1/p_X(x))\rceil\) satisfying
\begin{equation}
    H(X) \le \E{\ell(X)} < H(X)+1.
\end{equation}
\end{theorem}

\begin{proof}
The upper bound follows from $\lceil u\rceil<u+1$:
\[
\E{\ell(X)} = \sum_x p_X(x)\,\lceil \log(1/p_X(x))\rceil < \sum_x p_X(x)\bigl(\log(1/p_X(x))+1\bigr)=H(X)+1.
\]
The lower bound is Theorem~\ref{thm:entropy-lb}.
\end{proof}

\paragraph{Dyadic sources.}
If $p_X(x)$ are exact powers of two (``dyadic''), then $\log(1/p_X(x))$ are integers and the Shannon code achieves $\E{\ell(X)}=H(X)$.



\subsection{Cost of Mismatch in Compression}
When the true source distribution is $P_X$ but an estimated distribution $Q$ is used to design the code, the codeword lengths are assigned as
\[
l(x)=\log\frac{1}{q(x)}.
\]
Then, the expected length under the true distribution is
\[
\mathbb{E}_{P_X}[l(x)] = \sum_{x} p_X(x)\log\frac{1}{q_X(x)}.
\]
We can decompose this as
\[
\mathbb{E}_{P_X}[l(x)] = \sum_{x} p_X(x)\log\frac{p_X(x)}{q(x)} + \sum_{x} p_X(x)\log\frac{1}{p_X(x)} = D(P_X\|Q) + H(X).
\]
Thus, the additional cost due to the mismatch is exactly the KL divergence $D(P_X\|Q)$.


\subsection{Huffman Coding}
Huffman coding is an algorithm that produces an optimal prefix-free code for a given source distribution.

\subsubsection*{Huffman Algorithm}
\begin{enumerate}
    \item Identify the two symbols with the smallest probabilities.
    \item Merge these two symbols into a composite symbol whose probability is the sum of the two.
    \item Repeat the process on the reduced set of symbols until only one symbol remains, thereby forming a binary tree.
\end{enumerate}

\subsubsection*{Optimality Conditions}
Huffman codes satisfy the following conditions:
\begin{itemize}
    \item If $p(x_1) > p(x_2)$, then the optimal code must have $\ell(x_1) \le \ell(x_2)$.
    \item The two least probable symbols $a_{M-1}$ and $a_M$ receive codewords of equal length.
\end{itemize}

\begin{proof}(Sketch for condition 1)
Assume, to the contrary, that there exist symbols $x_1$ and $x_2$ with $p(x_1) > p(x_2)$ but $\ell(x_1) > \ell(x_2)$. By swapping the codewords of $x_1$ and $x_2$, the overall expected codeword length decreases, which contradicts the assumption of optimality.\\\\
(Sketch for condition 2) Suppose $l_M > l_{M-1}$. We can immediately find a better code since the code has a binary tree representation.
\end{proof}

\begin{figure}[h]
    \centering
    \includegraphics[width=0.7\textwidth]{Huffman.png}
    \caption{Example of a Huffman coding tree.}
    \label{fig:huffman_example}
\end{figure}

Although Huffman coding minimizes the expected codeword length among symbol-by-symbol prefix-free codes, we have
\[
H(X) \le \mathbb{E}[\ell(X)] < H(X) + 1,
\]
with the gap due to the integer constraint on the codeword lengths. This gap can be reduced by encoding blocks of symbols.


\subsection{Typicality and the Asymptotic Equipartition Property}
\label{sec:aep}

For block coding, the key concept is \emph{typical sequences}.
Let $X^n=(X_1,\dots,X_n)$ be i.i.d. with pmf $p_X$.
The joint pmf is $p(x^n)=\prod_{i=1}^n p_X(x_i)$.

\begin{definition}[$\varepsilon$-typical sequence]
Fix $\varepsilon>0$.
A sequence $x^n$ is \emph{$\varepsilon$-typical} if
\begin{equation}
    \left| -\frac{1}{n}\log p(x^n) - H(X) \right| \le \varepsilon.
\end{equation}
Let $A_\varepsilon^{(n)}$ be the set of all $\varepsilon$-typical sequences (the \emph{typical set}).
\end{definition}

\begin{theorem}[Asymptotic equipartition property (AEP)]
\label{thm:aep}
For any $\varepsilon>0$,
\begin{equation}
    \Pr{X^n\in A_\varepsilon^{(n)}} \xrightarrow[n\to\infty]{} 1.
\end{equation}
\end{theorem}

\begin{proof}
Let $Z_i\coloneqq -\log p_X(X_i)$.
Then $Z_1,Z_2,\dots$ are i.i.d. with $\E{Z_1}=H(X)$.
Moreover
\(-\frac{1}{n}\log p(X^n)=\frac{1}{n}\sum_{i=1}^n Z_i\).
By the weak law of large numbers, $\frac{1}{n}\sum_{i=1}^n Z_i\to H(X)$ in probability.
This is exactly the statement.
\end{proof}

\begin{theorem}[Size of the typical set]
\label{thm:typical-set-size}
For $n$ sufficiently large (depending on $\varepsilon$),
\begin{equation}
    (1-\varepsilon)\,2^{n(H(X)-\varepsilon)}
    \le |A_\varepsilon^{(n)}|
    \le 2^{n(H(X)+\varepsilon)}.
\end{equation}
\end{theorem}

\begin{proof}
If $x^n\in A_\varepsilon^{(n)}$, then
\(2^{-n(H+\varepsilon)}\le p(x^n)\le 2^{-n(H-\varepsilon)}\) where $H\equiv H(X)$.
For the upper bound,
\[
1\ge \Pr{X^n\in A_\varepsilon^{(n)}}=\sum_{x^n\in A_\varepsilon^{(n)}} p(x^n)
\ge |A_\varepsilon^{(n)}|\,2^{-n(H+\varepsilon)}.
\]
For the lower bound, Theorem~\ref{thm:aep} implies that for large $n$,
$\Pr{X^n\in A_\varepsilon^{(n)}}\ge 1-\varepsilon$.
Thus
\[
1-\varepsilon
\le \sum_{x^n\in A_\varepsilon^{(n)}} p(x^n)
\le |A_\varepsilon^{(n)}|\,2^{-n(H-\varepsilon)}.
\]
Rearrange both inequalities.
\end{proof}

\subsection{Shannon's Source Coding Theorem (Operational Form)}
\label{sec:shannon-source-coding}


Block coding allows us to approach the entropy limit arbitrarily closely.

\begin{theorem}[Achievability]
For any rate $R>H(X)+\epsilon$, there exists a compression scheme that maps the typical set $A_{\epsilon}^{(n)}$ into binary strings of length $\lceil nR\rceil$, ensuring that the probability of decoding error tends to zero as $n\to\infty$.
\end{theorem}
\begin{proof}
Fix $R > H(X) + \delta$. For $n$ sufficiently large, by the theorem 16 and 17, 
$$
\text{Pr}\left\{X^n \notin A_\delta^{(n)}\right\} < \epsilon
$$
and
$$
\left|A_\delta^{(n)}\right| \leq2^{n(H(X)+\delta)}=2^{nR}.
$$
Consider a scheme that enumerates sequences in $A_{\delta}^{(n)}$. The encoder maps $X^n \in A_\epsilon^{(n)}$ to its binary representation. Otherwise, it outputs $(0,0,\,...,\,0)$. The decoder maps this binary representation back to the corresponding sequence in $A_\delta^{(n)}.$ Thus, the error probability is bounded by
$$
P_e \leq \text{Pr}\left\{X^n \notin A_\delta^{(n)}\right\} < \epsilon
$$
and the rate is equal to
$$
\frac{\log \left|A_\delta^{(n)}\right|}{n} \leq \frac{nR}{n}=R.
$$
Thus, $R$ is an achievable rate.
\end{proof}

\begin{theorem}[Converse]
Any compression scheme with rate $R<H(X)$ cannot achieve near-lossless recovery as $n\to\infty$.
\end{theorem}
\begin{proof}
For $\delta > 0$, suppose $R < H(X)-\delta$. The size of possible output is $2^{nR}.$ Let $B$ denote the set of (uniquely decodable) reconstruction sequences $\hat{X}^n$. Then 
$$
\begin{aligned}
    \text{Pr}\left\{X^n \in B^c\right\} &\geq \text{Pr}\left\{X^n \in B^c \cap A_\delta^{(n)}\right\}\\
    &=\sum_{x^n\in B^c \cap A_\delta^{(n)}}p(x^n)\\
    &\geq \left|B^c \cap A_\delta^{(n)}\right|\cdot2^{-n(H(X)+\delta)}\\
    &\geq \left(|A_\delta^{(n)}|-\left|B\right|\right)\cdot2^{-n(H(X)+\delta)}\\
    & \approx (2^{nH(X)}-2^{nR})\cdot2^{-nH(X)}\\ &=1-2^{n(R-H(X))} \rightarrow 1 \text{ as } n\rightarrow \infty
\end{aligned}
$$
\end{proof}


\subsection{Arithmetic Coding}

In contrast to Huffman coding (which suffers from exponential complexity in block-length), arithmetic encoding efficiently represents sequences by mapping them to a subinterval of \([0,1)\).

\subsubsection*{Arithmetic Encoding}
Arithmetic encoding assigns a sequence \( x^n=(x_1,x_2,\ldots,x_n) \) an interval \([L,H)\) whose length equals the probability of the sequence:
\[
H - L = p(x^n) = p_X(x_1)\, p_X(x_2\,|\, x_1) \cdots p_X(x_n\,|\, x_1^{n-1}).
\]
Initially, the interval is \([0,1)\). This interval is recursively subdivided into subintervals according to the probabilities of the symbols in the alphabet. For example, if the alphabet is \(\{A, B, C\}\) with probabilities \(p_X(A)=p_A\), \(p_X(B)=p_B\), \(p_X(C)=p_C\), then the interval is partitioned proportionally to these probabilities.

\begin{example}[Encoding ``AB'']
Suppose the first symbol \(A\) maps to the interval \([0,0.25)\) and \(B\) maps to \([0.25,0.4)\) (according to their probabilities). Then, the string “AB” is encoded as the interval \([0.25,0.4)\). Choose 
\[
Z = \frac{L+H}{2} = 0.325,
\]
whose binary expansion is approximately
\[
0.325 \approx 0.0101001100\ldots_{(2)}.
\]
\end{example}

\subsubsection*{Truncation and Decoding}

Since the binary expansion of \(Z\) may be infinite, we choose to truncate it. Write
\[
Z = 0.Z_1Z_2\cdots Z_K\cdots, \quad \hat{Z} = 0.Z_1Z_2\cdots Z_K.
\]
We select the smallest integer \(K\) such that
\[
L \le \hat{Z} < \hat{Z} + 2^{-K} < H.
\]
This condition guarantees that the truncated value \(\hat{Z}\) still uniquely identifies a number within the interval \([L,H)\).


\begin{example}[Truncation for``AB'']
For the interval \([0.25,0.4)\) and \(Z=0.325\), if \(K=2\) then \(\hat{Z} = 0.01_{(2)} = 0.25\); however, \(\hat{Z} + 2^{-2} = 0.25+0.25 = 0.5\) is not contained in \([0.25,0.4)\). With \(K=3\), \(\hat{Z} = 0.010_{(2)} = 0.25\) and \(\hat{Z} + 2^{-3} = 0.25+0.125=0.375\), which lies in \([0.25,0.4)\). Hence, \(K=3\) is sufficient.
\end{example}

Decoding proceeds by using the received truncated binary number \(\hat{Z}\) to reconstruct the interval partitioning—that is, the decoder repeats the subdivision process to identify the original sequence \(x^n\).

\subsubsection*{Properties and Optimality of Arithmetic Encoding}
Let us prove the key property of arithmetic encoding.

\begin{theorem}[Property of Arithmetic Encoding]
Let \(Z \in [L,H)\) be represented by a truncated binary expansion using \(k\) bits such that
\[
L \le \hat{Z} < \hat{Z} + 2^{-k} < H.
\]
Then,
\[
2^{-k} \ge \frac{1}{4}(H-L).
\]
\end{theorem}
\begin{proof}
Assume the truncated value is 
\[
\hat{Z} = 0.x_1x_2\cdots x_k,
\]
and let \(\tilde{Z}\) denote the truncation to \(k-1\) bits:
\[
\tilde{Z} = 0.x_1x_2\cdots x_{k-1}.
\]
By the minimality of \(k\) (i.e., \(k\) is chosen as the smallest integer for which the condition holds), one of the following must occur:
\[
\tilde{Z} < L \quad \text{or} \quad \tilde{Z} + 2^{-(k-1)} > H.
\]
Suppose, without loss of generality, that
\[
\tilde{Z} < L.
\]
Then, since the condition fails for \(k-1\) bits, we must have 
\[
\tilde{Z} + 2^{-(k-1)} > \frac{L+H}{2}.
\]
Rearranging yields
\[
2^{-(k-1)} > \frac{H-L}{2},
\]
which implies
\[
2^{-k} \ge \frac{1}{4}(H-L).
\]
A similar argument holds if \(\tilde{Z} + 2^{-(k-1)} > H\).
\end{proof}

Using this inequality, one can show that if a sequence \(x^n\) is encoded by arithmetic encoding, then the codeword length \(\ell(x^n)\) satisfies
\[
\ell(x^n) \le \log \frac{1}{p(x^n)} + 2.
\]
Hence, the average length obeys
\[
\mathbb{E}[\ell(X^n)] \le H(X^n) + 2.
\]
Dividing by \(n\), the rate satisfies
\[
R \le H(X) + \frac{2}{n},
\]
so that for large \(n\) the rate approaches the entropy—demonstrating near-optimality.










\subsection{Sources with Memory and Entropy Rate} \label{sec:entropy-rate}
\subsubsection*{First-Order Markov Process}
\begin{definition}
A sequence \(X_1, X_2, \dots, X_n\) is said to follow a \emph{first-order Markov process} if
\[
p_X(x_i\,|\, x_1,x_2,\dots,x_{i-1}) = p_X(x_i\,|\, x_{i-1})
\]
for all \(i \ge 2\).
\end{definition}
Then, the joint distribution factors as
\[
p(x^n) = p_X(x_1) \prod_{i=2}^{n} p_X(x_i \,|\, x_{i-1}),
\]
and the joint entropy is
\[
H(X^n) = H(X_1) + \sum_{i=2}^{n} H(X_i \,|\, X_{i-1}).
\]
If the chain is \emph{stationary} (i.e., the marginal distribution of each \(X_i\) is the same), then
\[
H(X^n) = H(X_1) + (n-1)H(X_2 \,|\, X_1),
\]
so that the fundamental limit per symbol is
\[
\lim_{n \to \infty} \frac{1}{n} H(X^n) = H(X_2 \,|\, X_1).
\]

\begin{example}[Random Walk]
Consider the random walk defined by
\[
X_0 = 0,\quad X_n = X_{n-1} \pm 1 \quad \text{with probability } \frac{1}{2} \text{ each}.
\]
This process is Markov (i.e., the future depends only on the present). (Note: A simple random walk is not stationary since its variance grows with \(n\); here the example serves to illustrate the Markov property.)
\end{example}

\subsubsection*{\(k\)th-Order Markov Process}
\begin{definition}
A sequence \(X_1, X_2, \dots, X_n\) follows a \emph{\(k\)th-order Markov process} if
\[
p_X(x_i\,|\, x_1,\dots,x_{i-1}) = p_X(x_i\,|\, x_{i-k},\dots,x_{i-1})
\]
for all \(i > k\).
\end{definition}
Then,
\[
p(x^n) = \left[\prod_{i=1}^{k} p_X(x_i)\right] \prod_{i=k+1}^{n} p_X(x_i \,|\, x_{i-k}^{i-1}),
\]
and the fundamental limit becomes
\[
\lim_{n \to \infty}\frac{1}{n}H(X^n) = H(X_{k+1}\,|\, X_1^k).
\]

\subsubsection*{Stationary Distribution and Entropy Rate}
\begin{definition}[Entropy Rate]
For a stationary process \(\{X_i\}_{i=1}^{\infty}\), the \emph{entropy rate} is defined as
\[
H(\mathscr{X}) = \lim_{n\to\infty} \frac{1}{n} H(X_1,X_2,\dots,X_n).
\]
\end{definition}
For the i.i.d. case, \(H(\mathscr{X})=H(X)\); for a stationary first-order Markov process,
\[
H(\mathscr{X})=H(X_2\,|\, X_1);
\]
and for a stationary \(k\)th-order Markov process,
\[
H(\mathscr{X})=H(X_{k+1}\,|\, X^k).
\]

\begin{example}[Stationary Firsto-Order Markov]
Assume a process over three symbols \(\{A,B,C\}\) with an initial uniform distribution, i.e., \( p_X(A)=p_X(B)=p_X(C)=\frac{1}{3} \), and transition probabilities
\[
p_X(x_2\,|\, x_1) = \frac{1}{3} \quad \text{for all } x_1,x_2.
\]
Then the process is stationary and
\[
H(\mathscr{X}) = H(X_2\,|\, X_1) = \log 3.
\]
\end{example}

\begin{figure}[htpb]
\begin{center}
\includegraphics[scale=0.3]{scribe12-markovex.png}
\caption{\(1^\text{st}\) order Markov transition example.}
\label{fig:markovtrans}
\end{center}
\end{figure}

\subsubsection*{Achievability of the Fundamental Limit}
Consider a stationary first-order Markov source \(X_1,X_2,\ldots,X_n\). By the Markov property,
\[
p(x^n)=p_X(x_1)\prod_{i=2}^{n} p_X(x_i\,|\, x_{i-1}).
\]
Arithmetic encoding maps \(x^n\) to an interval of length
\[
H - L = p(x^n).
\]
Then, as shown earlier, the expected codeword length satisfies
\[
\mathbb{E}[\ell(X^n)] \le \log\frac{1}{p(x^n)} + 2.
\]
Taking the expectation and normalizing by \(n\), and under the stationarity assumption, we obtain
\[
R \leq \frac{1}{n}\left( H(X_1) + 2 \right) + \frac{n-1}{n}H(X_2\,|\, X_1).
\]
Letting \(n\to\infty\) yields
\[
R \approx H(X_2\,|\, X_1) = H(\mathscr{X}),
\]
establishing the achievability of the fundamental limit. A similar derivation holds for stationary \(k\)th-order Markov sources, for which the limit is \( H(\mathscr{X})=H(X_{k+1}\,|\, X^k) \).

\begin{remark}
In these notes we restrict our attention to stationary sources. The fundamental limit as given by the entropy rate holds only under stationarity.
\end{remark}



\subsection{Universal Compression: Learning While Coding}

Assume \(X_1, \dots, X_m\) are i.i.d.\ random variables with \(X_i \sim P_X\), but the true distribution \(P_X\) is unknown. One universal compression scheme is to use the first \(m\) symbols both to learn the distribution and to “synchronize” the encoder and decoder, and then to compress the remaining \(X_{m+1},\dots,X_n\) using a code designed from the estimated distribution.

A simple scheme is as follows:
\begin{enumerate}
    \item Examine (and encode) \(X_1,\dots,X_m\) using a pre-agreed, fixed (sub-optimal) code.
    \item Use these \(m\) symbols to compute an empirical estimate \(\hat{P}_X\) of the true distribution \(P_X\).
    \item Both encoder and decoder construct a Huffman code (or another optimal prefix-free code) based on \(\hat{P}_X\).
    \item Compress the remaining sequence \(X_{m+1},\dots,X_n\) using the newly constructed code.
\end{enumerate}

There is an inherent trade-off in the choice of \(m\): a larger \(m\) improves the accuracy of \(\hat{P}_X\) (and thus the compression rate may approach \(H(X)\)), but more bits are wasted to transmit the first \(m\) symbols using the fixed code. Conversely, a smaller \(m\) reduces the learning overhead but may yield a poor estimate of \(P_X\).

\subsubsection*{Lempel-Ziv Algorithms} \label{sec:lz}
Unlike the scheme above, which separates the learning and execution phases, Lempel-Ziv algorithms perform both simultaneously.
The LZ78 algorithm parses an input sequence into variable-length phrases while dynamically building a dictionary. The parsing rule is:
\begin{enumerate}
    \item Initialize the dictionary to be empty.
    \item Starting at the beginning of the sequence, find the longest prefix \(w\) that is already in the dictionary.
    \item Output the tuple \((d, a)\), where \(d\) is the index of \(w\) in the dictionary (or 0 if \(w\) is empty) and \(a\) is the next symbol that extends \(w\) (thus \(w a\) is the new phrase).
    \item Add \(w a\) to the dictionary.
    \item Continue parsing from the symbol following \(w a\).
\end{enumerate}

\begin{example}
Consider a sequence over the alphabet \(\{A, B\}\):
\[
\texttt{ABBABABAAABB}\ldots
\]
The parsing (with dictionary updates) proceeds as follows:
\begin{itemize}
    \item \textbf{Step 1:} The dictionary is initially empty. The first symbol \(A\) is not in the dictionary. Output \((0, A)\) and add \(A\) as phrase 1.
    \item \textbf{Step 2:} Next, consider \(B\). It is new, so output \((0, B)\) and add \(B\) as phrase 2.
    \item \textbf{Step 3:} The next symbols are \(B\). The algorithm finds that \(B\) is already in the dictionary (as phrase 2). Now, check the extension: after \(B\) comes \(A\) (i.e. the string “BA” is not in the dictionary). Output \((2, A)\) and add “BA” as phrase 3.
    \item \textbf{Step 4:} Continue with the next unparsed symbol. The sequence now starts with \(B\). \(B\) is in the dictionary. Check the extension: \(B\) followed by \(A\) gives “BA,” which is in the dictionary (phrase 3); then extend further: “BA” followed by \(B\) is “BAB,” which is new. Output \((3, B)\) and add “BAB” as phrase 4.
    \item \textbf{Step 5:} Next, the remaining portion begins with \(A\). Since \(A\) is already in the dictionary, check extension. If “A” followed by \(A\) (i.e. “AA”) is not in the dictionary, output \((1, A)\) and add “AA” as phrase 5.
    \item \textbf{Step 6:} Continue in this fashion for the rest of the input.
\end{itemize}
\end{example}

Each phrase is then encoded using an efficient binary representation of its tuple. Over time, frequently occurring phrases lead to long dictionary entries.

\begin{figure}[h]
    \centering
    \includegraphics[width=0.8\textwidth]{scribe14_figure.png}
    \caption{Mapping input symbols into a sequence of tuples (phrases) in LZ78.}
    \label{fig:lz78}
\end{figure}

Lempel-Ziv algorithms, like LZ78, are universal, meaning they asymptotically achieve the entropy rate \(H(\mathscr{X})\) for any stationary source \(\mathbb{X}\) without prior knowledge of \(P_X\).



